\section{Giới thiệu}

\begin{frame}{Công việc thực hiện}
  \begin{itemize}
    \item \textbf{Thách thức trong dự báo chất lượng nước biển:}
      \begin{itemize}
        \item Các trạm quan trắc (sensor) giới hạn về số lượng
        \item Cần đầy đủ các chỉ số quan trắc để đánh giá chất lượng nước biển
      \end{itemize}
    \item \textbf{Bài toán nội suy \& phóng chiếu:}
      \begin{itemize}
        \item \textbf{Phase 1 (Local Imputation):} Lấp đầy dữ liệu bị thiếu hụt ngay tại vị trí đặt trạm (suy diễn từ các chỉ số cảm biến cơ bản).
        \item \textbf{Phase 2 (Spatial Extrapolation):} Phóng chiếu dữ liệu xuyên không gian từ một Trạm Gốc (Anchor) sang vùng biển không có trạm (Target).
      \end{itemize}
    \item \textbf{Khuyết điểm của mô hình tuyến tính (MICE):} Khó có khả năng nắm bắt được quy luật phi tuyến do thời tiết và vi khí hậu tác động.
  \end{itemize}
\end{frame}
